% ID: SEL-D0060-OA-E03
% ESTADO: cotejada. Segunda pasada visual y borrador compilado revisados.
% Fuente: Selectividad/sel_2011_matematicas_aplicadas/223A-EXAMEN 3.pdf
% SHA-256: 5998699e0b2992b988a8e0ed0d2e3581dd8fac52d6f9802587ebd40fee5ef608
% Página: 1; opción A; ejercicio original 3.

En una ciudad, el $55\%$ de la población consume aceite de oliva, el $30\%$ de girasol, y el $20\%$ ambos tipos de aceite. Se escoge una persona al azar:
\begin{enumerate}[label=\alph*)]
\item \textbf{(1 punto)} Si consume aceite de oliva, \textquestiondown{}cuál es la probabilidad de que consuma también aceite de girasol?
\item \textbf{(1 punto)} Si consume aceite de girasol, \textquestiondown{}cuál es la probabilidad de que no consuma aceite de oliva?
\item \textbf{(0.5 puntos)} \textquestiondown{}Cuál es la probabilidad de que no consuma ninguno de los dos tipos de aceite?
\end{enumerate}
